\section{Introduction}
\label{sec:intro}

The computational neuroscience and machine learning communities have developed a multitude of methods to quantify similarity in population-level activity patterns across neural systems.
Indeed, a recent review by \textcite{klabunde2023similarity} catalogues over thirty approaches to quantifying representational similarity.
Many of these measures quantify similarity in the \textit{shape} or \textit{representational geometry} of point clouds.
For example, recent papers (e.g. \cite{williams2021generalized,Ding2021}) leverage the Procrustes distance and other concepts from \textit{shape theory}---an established body of work that formalizes the notion of shape and ways to measure distance between shapes~\cite{kendall2009shape,dryden2016statistical}.
Other measures of representational similarity are not quite this explicit but still emphasize desired \textit{geometric invariance} properties.
For example, work by \textcite{kornblith2019} popularized centered kernel alignment (CKA) by emphasizing its invariance to isotropic scaling, translation, and orthogonal transformations.
These are precisely the same invariances specified by classical shape theory~\cite{kendall2009shape,dryden2016statistical}.
Earlier work by \textcite{raghu2017svcca} argued for a more flexible invariance to affine transformations, and they used canonical correlations analysis (CCA) for this purpose.
Contemporary work in neuroscience is replete with similar geometric reasoning and quantitative frameworks~\cite{Kriegeskorte2021}.

% Here we ask: what does the \textit{shape} or \textit{geometry} of a neural representation imply about the \textit{function} performed by that neural system?
Here we ask: What does the similarity between neural representations, potentially in terms of shape or geometry, imply about the similarity between \emph{functions} performed by those neural systems?
Potentially, not very much.
\textcite{Maheswaranathan2019} showed that recurrent neural networks with different architectures performed the same task with similar dynamical algorithms, but with different representational geometry.
More recently, \textcite{lampinen2024learned} showed that representational geometry can be biased by other nuisance variables, such as the order in which multiple tasks are learned.
Similar limitations of representational geometry measures are discussed in~\cite{dujmovic2022some,davari2023reliability}.

On the other hand, several basic observations suggest that geometry and function ought to be related.
Consider, for example, the common practice of using linear models to decode task-relevant variables from neural population activity~\cite{alain2017understanding,Kriegeskorte2019}.
The premise behind these analyses is that anything linearly decodable from system $A$ is readily accessible to layers or brain regions immediately downstream of $A$.
Therefore, information that is linearly decodable from $A$ may relate to functional role played by $A$ within the overall system \cite{Cao2022}.
Notably, the invariances of typical representational similarity measures---translations, isotropic scalings, and orthogonal transformations---closely coincide with the set of transformations that leave decoding accuracy unchanged.
For example, translations and rotations of neural population activity will not impact the accuracy of a linear decoder with an intercept term and an $L_2$ penalty on the weights.
Thus, if we accept the premise that decoding accuracy is a proxy---even, perhaps, an imperfect proxy---for neural system function, then representational geometry may be a reasonable framework to capture something about this function~\cite{kriegeskorte2019peeling}.

We provide a unifying theoretical framework that connects several existing methods of measuring representational similarity in terms of linear decoding statistics.  
Specifically, we show that popular methods such as centered kernel alignment (CKA) \cite{kornblith2019} and similarity based on canonical correlations analysis (CCA) \cite{raghu2017svcca} can be interpreted as alignment scores between optimal linear decoders with particular weight regularizations.  
Additionally, we study how the \emph{shape} of neural representations is related to decoding by deriving bounds that relate the average decoding distance and the Procrustes distance.  
We find that the Procrustes distance is a more 
strict notion of geometric dissimilarity than the expected distance between optimal decoder readouts, in that a small value of the Procrustes distance implies a small expected distance between optimal decoder readouts but the converse is not necessarily true.  
This formalizes recent observations by \textcite{cloos2024differentiable}, who found that high Procrustes similarity implied high CKA similarity in practical settings.

